\section{Private norm primes and higher moments on a positive-density class}
\label{pn-section}

For a positive integer \(B\), retain the actual roots
\[
 a_k=3k,\quad w_k=a_k+\zeta,\quad Q_k=a_k^2+a_k+1,\quad
 \alpha_k=w_k/\bar w_k,\qquad B\leq k<2B.
\]
Define the actual subset
\[
 \mathcal{S}_B=\{k:B\leq k<2B,\ Q_k
                    \text{ has a prime divisor greater than }12B\}.
\]
This section proves membership in a useful subclass; it does not
assume that the actual roots have independent random prime factors.

\begin{theorem}[A positive-density class of independent ratios]
\label{pn-private}
As \(B\to\infty\),
\[
 |\mathcal{S}_B|\geq\frac B2
               -O\!\left(\frac{B\log\log B}{\log B}\right).
\]
For every \(k\in\mathcal{S}_B\) there is exactly one prime
\(r_k>12B\) in \(Q_k\), its depth is one, and it divides no other
norm in the entire block. The algebraic numbers
\(\{\alpha_k:k\in\mathcal{S}_B\}\) are multiplicatively independent.
\end{theorem}
\begin{proof}
Uniformly in the block,
\[
 9B^2\leq Q_k<49B^2,\qquad
 \sum_k\log Q_k=2B\log B+O(B).
\]
The polynomial \(9X^2+3X+1\) has no factor two or three.
Every supported prime is \(1\pmod3\), and at such a prime \(r\)
it has exactly two simple roots, each with one lift at every
depth. Its discriminant is \(-27\).
Put \(H_r=\lfloor\log(49B^2)/\log r\rfloor\).
Counting the actual interval at each depth gives
\[
 \sum_k v_r(Q_k)\log r
     \leq\frac{2B\log r}{r-1}+2H_r\log r.
\]

We use only the fixed-modulus consequence
\[
 \vartheta(x;3,1)=\frac x2+O(x/\log x)
\]
of Theorem~1.2, formula~(1.12), of
\href{https://arxiv.org/abs/1802.00085v3}
{Bennett, Martin, O'Bryant and Rechnitzer}.
Partial summation and the convergent change from \(1/r\) to
\(1/(r-1)\) give
\[
 \sum_{\substack{r\leq x\\r\equiv1\ (3)}}\frac{\log r}{r-1}
       =\frac12\log x+O(\log\log x).
\]
Indeed the error integral is \(O(\int^x dt/(t\log t))\).
The sum of the endpoint terms at \(r\leq12B\) is \(O(B)\):
each is at most \(2\log(49B^2)\), and
\(\pi(12B)=O(B/\log B)\). Thus all small norm primes together
contribute at most
\[
 B\log B+O(B\log\log B).
\]
The mass above \(12B\) is at least
\(B\log B-O(B\log\log B)\).
Each root contributes at most \(2\log B+O(1)\), proving the count.
This is a fixed progression estimate, not a moving-field input.

Since \((12B)^2>49B^2\), a norm has at most one such prime,
and its valuation is one. If \(r>12B\) divided norms at two
distinct parameters \(a,b\), it would divide
\[
 Q(a)-Q(b)=(a-b)(a+b+1).
\]
Both nonzero factors have absolute value less than \(r\), which
is impossible. Finally, the primitive unramified \(w_k\) has
one orientation above \(r_k\). At that prime ideal \(\alpha_k\)
has valuation \(1\) or \(-1\), while every other block ratio
has valuation zero. Applying these valuations to a multiplicative
relation forces every exponent to vanish.
\end{proof}

\begin{theorem}[Higher moments at deep boundary primes]\label{pn-moment}
Fix an integer \(\nu\geq2\). For every sufficiently large prime
\(n\), put \(B=n^4\). For every prime \(q>B\),
\[
 \#\{k\in\mathcal{S}_B:q^{2\nu}\mid T_n(k)\}
             \leq(3\nu!n)^{1/\nu}.
\]
The sufficiently-large condition depends on the fixed \(\nu\).
\end{theorem}
\begin{proof}
At every hit, the root ratio lies in the norm-one group killed by
\(3n\) modulo \(q^{2\nu}\).
The split and inert lifting proof of Lemma~\ref{mc-torsion}
works at this precision too and bounds that group by \(3n\).

For two ordered \(\nu\)-tuples write
\[
 z=\prod_{j=1}^{\nu}(a_j+\zeta)=R+S\zeta,\qquad
 z'=R'+S'\zeta.
\]
The norm coordinate bound and a telescoping expansion give
\[
 |R|\leq2(7B)^\nu,\qquad
 \left|z-\prod_j a_j\right|\leq\nu(7B)^{\nu-1},
 \qquad |S|\leq2\nu(7B)^{\nu-1}.
\]
The middle comparison is with a real integer. Hence
\[
 |RS'-R'S|\leq C_\nu B^{2\nu-1},\qquad
 C_\nu=8\nu\,7^{2\nu-1}.
\]
When \(B>C_\nu\), this is smaller than \(q^{2\nu}\).
Equal tuple images imply divisibility of this integer by
\(q^{2\nu}\), so their actual algebraic ratio products are equal.
The independence in Theorem~\ref{pn-private} implies equal
multiplicity of every root index. Thus each image has at most
\(\nu!\) ordered tuples, also when entries repeat.
If the set of hits has size \(M\), its \(M^\nu\) tuples have at
most \(3n\) images. Therefore \(M^\nu\leq3n\nu!\).
No assertion uniform in a varying \(\nu\) is used.
\end{proof}

\begin{theorem}[Full positive excess in a farther interval]
\label{pn-excess}
Fix an integer \(\nu\geq3\) and \(0<\delta<1-1/\nu\), and for
sufficiently large prime \(n\) put
\[
 B=n^4,\qquad Z=\lfloor Bn^{1-1/\nu-\delta}\rfloor>B,\qquad
 F_{B,Z}(k)=\sum_{B<q\leq Z}(v_q(T_n(k))-3)_+\log q.
\]
Then
\[
 \frac1B\sum_{k\in\mathcal{S}_B}\frac{F_{B,Z}(k)}{t_k}
       =O_{\nu,\delta}\!\left(\frac{n^{-\delta}}{\log n}\right).
\]
\end{theorem}
\begin{proof}
Use the exact identity
\((e-3)_+=\sum_{j\geq4}{\bf1}_{e\geq j}\).
For each prime, every layer \(4\leq j<2\nu\) has at most
\(\sqrt{6n}\) roots in \(\mathcal{S}_B\), by
Theorem~\ref{pn-moment} with parameter two. Every higher layer
has at most \((3\nu!n)^{1/\nu}\) roots.
The elementary actual height bound caps the depth by
\[
 H_q\leq\frac{3nL}{\log q},\qquad L=\log(6B+1).
\]
Consequently the complete positive excess at a single prime obeys
\begin{align}
 \sum_{k\in\mathcal{S}_B}(v_q(T_n(k))-3)_+\log q
 &\leq(2\nu-4)\sqrt{6n}\log q \notag\\
 &\quad+3nL(3\nu!n)^{1/\nu}. \label{pn-layer-bill}
\end{align}
This remains an upper bound if a range of layers is empty.
It explicitly pays for the intermediate depths as well as the
deepest ones.

Since \(t_k\geq n\log(3B)\), \(q\leq Z\leq Bn\), and \(B=n^4\),
division by \(Bn\log(3B)\) bounds \eqref{pn-layer-bill} by
\(C_\nu n^{1/\nu}/B\).
A prime contributing positive excess has rank \(n\): rank one
would put \(q^2\) in one of \(a_k,a_k+1\), each below
\(6B+1<q^2\). Thus \(q\equiv\pm1\pmod{3n}\).
Brun--Titchmarsh bounds the number of eligible primes by
\[
 \frac{2Z}{(n-1)\log(Z/(3n))}.
\]
The total is therefore
\[
 O_\nu\!\left(\frac{Zn^{1/\nu}}{Bn\log n}\right)
       =O_{\nu,\delta}\!\left(\frac{n^{-\delta}}{\log n}\right).
\]
The normalization is by the entire block size \(B\), not an
unstated replacement by \(|\mathcal{S}_B|\).
\end{proof}

\begin{corollary}[An actual positive-density domain]\label{pn-domain}
For fixed \(\nu,\delta\) as above, there is a set
\(\mathcal{P}_{n,\nu,\delta}\) of actual roots of relative size at
least \(1/2-o(1)\), lying in \(\mathcal{S}_B\), such that
\[
 \frac{F_{B,Z}(k)}{t_k}\leq n^{-\delta/2},\qquad
 \frac{L_B(T_n(k))}{t_k}\leq(\log n)^{-1/2}.
\]
Writing
\[
 W_Z(k)=\sum_{\substack{q>Z\\q\mid T_n(k)}}
                         (v_q(T_n(k))-3)\log q,
\]
one has on this set
\begin{align}
 \frac{J(T_n(k))}{t_k}
   &\leq\frac{W_Z(k)}{t_k}+(\log n)^{-1/2}+n^{-\delta/2},
       \label{pn-signed}\\
 \frac{\log\operatorname{rad}(T_n(k))}{t_k}
   &\geq1-\frac4{3n}
       -\frac{(\log n)^{-1/2}+n^{-\delta/2}}3
       -\frac{(W_Z(k))_+}{3t_k}. \label{pn-radical}
\end{align}
\end{corollary}
\begin{proof}
Markov's inequality in Theorem~\ref{pn-excess} removes
\(O_{\nu,\delta}(n^{-\delta/2}/\log n)\) of the whole block.
The full-mass mean of Theorem~\ref{fml-mean}, with the elementary
inputs of Section~\ref{ebi-section}, at \(B\) is \(O(1/\log n)\).
Markov at \((\log n)^{-1/2}\) removes \(O((\log n)^{-1/2})\).
Intersect with Theorem~\ref{pn-private}.
The signed cost up to \(B\) is at most \(L_B\), and that between
\(B\) and \(Z\) is at most \(F_{B,Z}\).
This proves \eqref{pn-signed}; the angular estimate
\(\log T_n/t_k\geq3-4/n\) proves \eqref{pn-radical}.
\end{proof}

The intermediate interval need not have negligible full mass.
Its negative credits can be retained in a sharper application.
For arbitrarily large \emph{fixed} \(\nu\), its upper exponent
\(1-1/\nu-\delta\) can approach one; this does not interchange
the \(\nu\) and \(n\) limits. The private-norm condition is proved
for at least half the roots asymptotically, while the remaining
roots and the unbounded farther signed cost are not controlled.
The independent ordinary reviews cover these conclusions.
No formalization of the progression estimate, the private-norm
density theorem, or the higher-moment analytic consequences is
claimed. A vanishing positive far cost is still required for
the restricted ABC consequence along a sequence of these roots.
